\documentclass[11pt]{article}

\usepackage{sectsty}
\usepackage{siunitx}
\usepackage{graphicx}
\usepackage{amsmath}
\usepackage{amsthm}
\usepackage{float}
\usepackage{outlines}
\usepackage{mathtools}
\usepackage[thinc]{esdiff}
\usepackage{amssymb}
\usepackage{outlines}
\usepackage{caption}
\usepackage{subcaption}
\let\oldsi\si
\renewcommand{\si}[1]{\oldsi[per-mode=reciprocal-positive-first]{#1}}
\usepackage{enumitem}
\newcommand{\degsym}{^{\circ}}
\newcommand{\Mod}[1]{\ (\mathrm{mod}\ #1)}
\usepackage{hyperref}
\hypersetup{
    colorlinks,
    citecolor=black,
    filecolor=black,
    linkcolor=black,
    urlcolor=black
}
\newcommand{\lb}{\\[8pt]}
\newenvironment*{cell}[1][]{\begin{tabular}[c]{@{}c@{}}}{\end{tabular}}
\newcommand{\img}[3]{\begin{center}
  \begin{figure}[H]
    \centering
    \includegraphics[width=#2\textwidth]{#1}
    \caption{#3}
    \label{fig:fig1}
  \end{figure}
\end{center}}
\newcommand{\doubleimg}[4]{\begin{center}
  \begin{figure}[H]
    \centering
    \begin{subfigure}{.45\textwidth}
      \centering
      \includegraphics[width=1\linewidth]{#1}
      \caption{#2}
      \label{fig:sub1}
    \end{subfigure}
    \begin{subfigure}{.45\textwidth}
      \centering
      \includegraphics[width=1\linewidth]{#3}
      \caption{#4}
      \label{fig:sub2}
    \end{subfigure}
  \end{figure}
\end{center}}



\title{Math AA HL at KCA - Chapter 14 Notes}
\author{Tim Bao 2023-2025}
\date{March 11th, 2024}

\begin{document}

\maketitle
\pagebreak
\tableofcontents
\pagebreak

\section{Differentiation by First Principles}

The derivative of a function is defined using limits as $$f'(x) = \lim\limits_{h \to 0}\frac{f(x + h) - f(x)}{h}$$

\pagebreak

\section{Differentiation Rules}

\begin{itemize}
  \item $\diff{}{x}x^n = nx^{n - 1}$
  \item $\diff{}{x}\sin x = \cos x$
  \item $\diff{}{x}\cos x = -\sin x$
  \item $\diff{}{x}\tan x = \dfrac{1}{\cos^2 x} = \sec^2x = 1 + \tan^2 x$
  \item $\diff{}{x}\csc x = -\csc x \cot x$
  \item $\diff{}{x}\sec x = \tan x \sec x$
  \item $\diff{}{x}\cot x = -\csc^2 x$
  \item $\diff{}{x}\arcsin x = \dfrac{1}{\sqrt{1 - x^2}}$
  \item $\diff{}{x}\arccos x = -\dfrac{1}{\sqrt{1 - x^2}} = -(\arcsin x)'$
  \item $\diff{}{x}\arctan x = \dfrac{1}{1 + x^2}$
  \item $\diff{}{x}e^x = e^x$
  \item $\diff{}{x}a^x = a^x\cdot \ln a$
  \item Product rule: $\diff{}{x}f(x)g(x) = f'(x)g(x) + f(x)g'(x)$
  \item Quotient rule: $\diff{}{x}\dfrac{f(x)}{g(x)} = \dfrac{f'(x)g(x) - f(x)g'(x)}{g(x)^2}$
  \item Chain rule: $\diff{}{x}f(g(x)) = f'(g(x)) \times g'(x)$
\end{itemize}

\vspace{0.5cm}

\noindent Basic differentiation rules are usually combined through the chain rule.

\pagebreak

\section{Equations of Tangents and Normals}

Given a function $f(x)$, the equation of the tangent line at $x = a$ is $$y = f'(a)(x - a) + f(a)$$\lb
The normal at $x = a$ is given by $$y = -\frac{1}{f'(a)}(x - a) + f(a)$$

\vspace{2cm}

\section{Increasing and Decreasing Functions over an Interval}

\begin{itemize}
  \item $f'(a) = 0 \iff$ $(a, f(a))$ is a \textit{stationary point}
  \item $f'(a) > 0 $ for $x\in (a, b)$, then $f(x)$ is \textit{increasing} in the interval $(a, b)$
  \item $f'(a) < 0 $ for $x\in (a, b)$, then $f(x)$ is \textit{decreasing} in the interval $(a, b)$
\end{itemize}

\pagebreak

\section{Second Order Derivatives}

The second-order derivative of $y = f(x)$ can be written as

$$f''(x) = \diff[2]{y}{x}$$
Interpreting the second-order derivative
\begin{itemize}
  \item $f''(a) = 0 \iff$ $(a, f(a))$ is a point of inflection, if the graph alters its concavity (concave to convex or vice versa) at that point
  \item $f''(a) > 0 \iff$ convex/concave up at $x = a$; for a turning point, it is then a \textit{minimum point}
  \item $f''(a) < 0 \iff$ concave down at $x = a$; for a turning point, it is then a \textit{maximum point}
\end{itemize}

\section{Sketching Gradient Functions}

\bgroup
\def\arraystretch{1.5}%
\begin{center}
  \begin{tabular}{c|c}
    On $f(x)$            & On $f'(x)$                   \\
    \hline
    Turning point        & Cuts the $x$-axis            \\ \hline
    Inflection point     & Turning point                \\ \hline
    Positive gradient    & Above $x$-axis               \\ \hline
    Negative gradient    & Below $x$-axis               \\ \hline
    Vertical asymptote   & Vertical asymptote           \\ \hline
    Horizontal asymptote & Horizontal asymptote $y = 0$ \\ \hline
  \end{tabular}
\end{center}
\egroup

\pagebreak

\section{Implicit Differentiation}

\end{document}